\documentclass[a4paper,12pt]{article}
\usepackage[utf8]{inputenc}
\usepackage[T1]{fontenc}
\usepackage[french]{babel}
\usepackage{geometry}
\geometry{
    a4paper,
    left=20mm,
    right=20mm,
    top=20mm,
    bottom=20mm
}
\usepackage{graphicx}

\begin{document}
\pagestyle{empty}
\Large

% ==================== PAGE 1 ====================

% --- SLIDE 8 ---
\begin{minipage}[t][0.45\textheight]{\textwidth}
\textbf{\huge SLIDE 8 : Les Correctifs Blindés}\\
\normalsize \textit{Temps: 0:00 - 1:10}\\

\large
\textit{(Transition : Prendre la parole après le Speaker 3 pour montrer la solution)}\\

\textbf{Le script :}\\
« Nous venons de voir à quel point une configuration par défaut pouvait nous exposer au danger. Alors, comment reprendre le contrôle et colmater ces brèches ? C'est l'heure de la contre-attaque. Sur cet écran, vous avez sous les yeux notre code Terraform "blindé", conçu pour créer une infrastructure vraiment résiliente.\\

Si vous regardez l'éditeur au centre, nous avons appliqué cinq correctifs fondamentaux :\\
- \textbf{L'accès public est strictement bloqué}.\\
- Nous avons forcé le \textbf{chiffrement AES256}.\\
- \textbf{L'accès SSH est restreint} à notre VPN (\texttt{123.45.67.89}).\\

Enfin, nous avons activé le \textbf{versioning} et une \textbf{journalisation complète}. Nous avons maintenant des caméras de surveillance sur chaque action. Notre forteresse logicielle est construite. »
\end{minipage}

\vspace*{\fill}
\begin{center}
\noindent\dotfill\textit{\small \ \ Couper ici (Milieu de la page A4)\ \ }\dotfill\\
\end{center}
\vspace*{\fill}

% --- SLIDE 9 ---
\begin{minipage}[t][0.45\textheight]{\textwidth}
\textbf{\huge SLIDE 9 : La Validation Trivy}\\
\normalsize \textit{Temps: 1:10 - 2:10}\\

\large
\textit{(Transition : Cliquez sur la slide suivante. Montrez l'écran du doigt)}\\

\textbf{Le script :}\\
« Mais en cybersécurité, la confiance n'exclut pas le contrôle. Il faut des preuves. Que se passe-t-il si nous repassons notre radar de balayage, Trivy, sur cette nouvelle version de l'infrastructure ?\\

Regardez le terminal au centre. Le résultat est sans appel, c'est un succès total : \textbf{9 tests passés, 0 échec}. L'analyse confirme que l'accès public, le chiffrement, les logs et la restriction SSH sont tous au vert.\\

Zéro vulnérabilité critique. Aucune faille ne passera entre les mailles du filet. Cela signifie que l'infrastructure que nous avons codée est \textbf{techniquement validée} par nos outils de scan. Elle est prête pour affronter la production en toute sécurité. »
\end{minipage}

\newpage
% ==================== PAGE 2 ====================

% --- SLIDE 10 ---
\begin{minipage}[t][0.45\textheight]{\textwidth}
\textbf{\huge SLIDE 10 : Conclusion}\\
\normalsize \textit{Temps: 2:10 - 3:00}\\

\large
\textit{(Transition : Cliquez sur la dernière slide. Regardez directement l'audience, ton plus posé)}\\

\textbf{Le script :}\\
« Pour conclure notre présentation, qu'est-ce que nous devons retenir aujourd'hui ?\\

L'Infrastructure as Code (IaC) est formidable : elle rime avec efficacité. Mais un grand pouvoir implique de grandes responsabilités... Déployer vite, c'est aussi le risque de déployer des failles très vite. \\
C'est pourquoi \textbf{la sécurité doit commencer dès la conception}.\\

Les scanners automatisés, comme Trivy, doivent faire partie de votre pipeline. C'est en vérifiant notre code avant de le déployer que nous garantissons la résilience de notre cloud en continu.\\

Gardez cette phrase en tête : \textit{« In Code We Trust, But We Verify »}. C'est ainsi que l'on construit un avenir maîtrisé.\\

Merci beaucoup pour votre attention. S'il vous reste des questions, nous sommes à votre entière disposition ! »
\end{minipage}

\vspace*{\fill}
\begin{center}
\noindent\dotfill\textit{\small \ \ Couper ici (Milieu de la page A4)\ \ }\dotfill\\
\end{center}
\vspace*{\fill}

% --- BLANK EXTRA ---
\begin{minipage}[t][0.45\textheight]{\textwidth}
\begin{center}
\vspace*{3cm}
\textit{\Huge (Espace vide pour annotations)}
\end{center}
\end{minipage}

\end{document}